\chapter{Odwzorowania liniowe, jądro i obraz}
\label{ch:lin3-odwzorowania-liniowe}
\index{odwzorowanie liniowe@Odwzorowanie liniowe}
\index{jadro odwzorowania@Jądro odwzorowania liniowego}
\index{obraz odwzorowania@Obraz odwzorowania liniowego}

\noindent\textbf{Cel rozdziału.}
W Rozdziale~\ref{ch:lin-przestrzenie-wektorowe}
poznaliśmy przestrzenie wektorowe, bazy Hamela i współrzędne,
a w Rozdziale~\ref{ch:lin2-wymiar-sumy-ilorazy}
sumy proste, wymiar kardynalny i przestrzenie ilorazowe.
Teraz badamy odwzorowania zachowujące strukturę liniową.
Wskażemy, jak ich jądro mierzy utratę informacji,
jak obraz opisuje osiągalne wartości i dlaczego
iloraz przez jądro jest izomorficzny z obrazem.
Wyznaczymy również przestrzeń wszystkich odwzorowań
liniowych i omówimy projekcje.

W całym rozdziale \(K\) jest ciałem, \(0_K\neq1_K\),
a wszystkie przestrzenie wektorowe są nad \(K\).
Nie zakładamy skończonej wymiarowości, o ile nie
zostanie ona wyraźnie wskazana. Suma wektorów
indeksowana nieskończonym zbiorem oznacza
wyłącznie sumę o skończonym nośniku.
Przy wynikach wymagających wyboru baz dowolnych
przestrzeni nieskończenie wymiarowych
jawnie przyjmujemy aksjomat wyboru.

\section{Definicja i pierwsze własności}
\label{sec:lin3-definicja}

\begin{definicja}[Odwzorowanie liniowe]
\label{def:lin3-odwzorowanie-liniowe}
Niech \(V,W\) będą przestrzeniami wektorowymi nad \(K\).
Odwzorowanie \(T\colon V\to W\) nazywamy
\emph{liniowym}, jeżeli dla dowolnych \(v,w\in V\)
i \(a\in K\) zachodzą
\[
 T(v+w)=T(v)+T(w),\qquad T(av)=aT(v).
\]
Pierwszy warunek nazywamy \emph{addytywnością},
a drugi \emph{jednorodnością}.
\end{definicja}

\begin{stwierdzenie}[Kryterium jednej kombinacji]
\label{prop:lin3-kryterium}
Niech \(V,W\) będą przestrzeniami wektorowymi
nad \(K\), a \(T\colon V\to W\) odwzorowaniem.
Jest ono liniowe wtedy i tylko wtedy, gdy
dla wszystkich \(a,b\in K\) oraz \(v,w\in V\) zachodzi
\[
 T(av+bw)=aT(v)+bT(w).
\]
\end{stwierdzenie}
\begin{proof}
Jeżeli \(T\) jest liniowe, to stosujemy
addytywność do \(av+bw\), a następnie jednorodność
do każdego składnika. W przeciwnym kierunku
ustawmy \(a=b=0_K\); otrzymujemy \(T(0_V)=0_W\).
Następnie dla \(a=b=1_K\) otrzymujemy
\(T(v+w)=T(v)+T(w)\).
Wreszcie przy \(b=0_K\) dostajemy
\(T(av)=aT(v)+0_KT(w)=aT(v)\).
Oba warunki definicji są więc spełnione.
\end{proof}

\begin{lemat}[Podstawowe tożsamości]
\label{lem:lin3-podstawowe}
Niech \(T\colon V\to W\) będzie liniowe,
gdzie \(V,W\) są przestrzeniami wektorowymi nad \(K\).
Wtedy \(T(0_V)=0_W\), \(T(-v)=-T(v)\) dla
każdego \(v\in V\), a dla dowolnych
\(m\in\mathbb N_0\), \(a_1,\ldots,a_m\in K\)
i \(v_1,\ldots,v_m\in V\) mamy
\[
 T\left(\sum_{j=1}^m a_jv_j\right)
   =\sum_{j=1}^m a_jT(v_j).
\]
Równość obejmuje przypadek \(m=0\).
\end{lemat}
\begin{proof}
Z jednorodności \(T(0_Kv)=0_KT(v)=0_W\).
Podstawiając skalar \(-1_K\), otrzymujemy
\(T(-v)=(-1_K)T(v)=-T(v)\).
Równość dla sumy jednego składnika wynika
z jednorodności. Jeżeli zachodzi dla \(m\),
to dla \(m+1\) używamy addytywności,
a następnie założenia indukcyjnego:
\[
 T\left(\sum_{j=1}^{m+1}a_jv_j\right)
 =T\left(\sum_{j=1}^{m}a_jv_j\right)
   +T(a_{m+1}v_{m+1})
 =\sum_{j=1}^{m+1}a_jT(v_j).
\]
Przy \(m=0\) obie strony są wektorami zerowymi.
\end{proof}

\begin{przyklad}[Odwzorowania współrzędnych]
\label{prz:lin3-wspolrzedne}
Niech \(K\) będzie dowolnym ciałem i oznaczmy
\(2_K:=1_K+1_K\).
Odwzorowanie
\[
 T\colon K^2\longrightarrow K^3,\qquad
 T(x,y)=(x+y,x-y,2_Kx)
\]
jest liniowe. Istotnie, w każdej współrzędnej
otrzymujemy liniową kombinację \(x,y\).
Dla \(u=(x,y)\), \(v=(s,t)\) i \(a,b\in K\)
pierwsza współrzędna \(T(au+bv)\) wynosi
\(a(x+y)+b(s+t)\); w dwóch pozostałych
rachunek jest analogiczny. Daje to
\(T(au+bv)=aT(u)+bT(v)\).
W przypadku charakterystyki \(2\)
trzecia współrzędna jest zawsze zerowa;
nie wpływa to na liniowość odwzorowania.
\end{przyklad}

\begin{przyklad}[Przykłady niebędące liniowymi]
\label{prz:lin3-nieliniowe}
Niech \(S\colon V\to W\) będzie liniowe
i niech \(w_0\in W\setminus\{0_W\}\).
Odwzorowanie afiniczne \(A(v)=S(v)+w_0\)
nie jest liniowe, gdyż \(A(0_V)=w_0\neq0_W\).
Także \(f\colon\mathbb R\to\mathbb R\),
\(f(x)=x^2\), nie jest liniowe:
\(f(1+1)=4\), natomiast \(f(1)+f(1)=2\).
Ostatnie równanie jest kontrprzykładem
do addytywności, a nie tylko do jednorodności.
\end{przyklad}

\section{Odwzorowanie liniowe wyznaczone na bazie}
\label{sec:lin3-baza}

\begin{twierdzenie}[Przedłużanie zadanych wartości na bazie]
\label{thm:lin3-przedluzenie-z-bazy}
Niech \(V,W\) będą przestrzeniami wektorowymi nad \(K\),
\(B=(b_i)_{i\in I}\) bazą Hamela \(V\), a
\((w_i)_{i\in I}\) dowolną rodziną wektorów \(W\).
Wtedy istnieje dokładnie jedno odwzorowanie liniowe
\(T\colon V\to W\) takie, że \(T(b_i)=w_i\)
dla każdego \(i\in I\). Określa je wzór
\[
 T\left(\sum_{i\in F}a_ib_i\right)
   =\sum_{i\in F}a_iw_i,
 \qquad F\subseteq I\text{ skończony},\quad a_i\in K.
\]
Jeżeli baza \(B\) jest już dana, twierdzenie
nie wymaga aksjomatu wyboru.
\end{twierdzenie}
\begin{proof}
Każdy \(v\in V\) ma dokładnie jedną rodzinę
współczynników w bazie \(B\), o skończonym nośniku,
na mocy Twierdzenia~\ref{thm:lin-wspolrzedne-hamela}.
Zatem prawa strona wzoru jest dobrze określona:
różne przedstawienia tego samego \(v\),
po dopisaniu zer, mają identyczne współczynniki.
Otrzymujemy więc odwzorowanie \(T\colon V\to W\).

Jeśli \(v=\sum_{i\in F}a_ib_i\) i
\(u=\sum_{i\in G}c_ib_i\), to dopisujemy zera
poza odpowiednimi nośnikami. Dla \(p,q\in K\)
współczynniki \(pv+qu\) względem \(B\)
wynoszą \(pa_i+qc_i\) dla \(i\in F\cup G\).
Stąd
\[
 \begin{aligned}
 T(pv+qu)
 &=\sum_{i\in F\cup G}(pa_i+qc_i)w_i\\
 &=p\sum_{i\in F}a_iw_i
   +q\sum_{i\in G}c_iw_i
 =pT(v)+qT(u).
 \end{aligned}
\]
Odwzorowanie jest liniowe i spełnia
\(T(b_i)=w_i\), gdyż współrzędne \(b_i\)
są równe \(1_K\) przy \(i\) i zerowe poza \(i\).

Jeśli \(S\colon V\to W\) jest innym odwzorowaniem
liniowym o tych samych wartościach na bazie,
to dla dowolnego \(v=\sum_{i\in F}a_ib_i\)
z Lematu~\ref{lem:lin3-podstawowe} otrzymujemy
\[
 S(v)=\sum_{i\in F}a_iS(b_i)
      =\sum_{i\in F}a_iw_i=T(v).
\]
Zatem \(S=T\). Przy \(I=\varnothing\)
przestrzeń \(V\) jest zerowa, a jedynym
odwzorowaniem liniowym \(V\to W\)
jest odwzorowanie zerowe; powyższy argument
obejmuje również ten przypadek.
\end{proof}

\begin{wniosek}[Zgodność na generatorach]
\label{cor:lin3-zgodnosc-na-generatorach}
Niech \(V,W\) będą przestrzeniami wektorowymi
nad \(K\), a \(S,T\colon V\to W\)
odwzorowaniami liniowymi. Jeżeli
\(A\subseteq V\) rozpina \(V\) i \(S(a)=T(a)\)
dla każdego \(a\in A\), to \(S=T\).
\end{wniosek}
\begin{proof}
Dla każdego \(v\in V\) istnieją
\(m\in\mathbb N_0\), skalary \(c_j\in K\)
i wektory \(a_j\in A\), \(1\leq j\leq m\),
takie, że \(v=\sum_{j=1}^{m}c_ja_j\).
Liniowość obu odwzorowań i ich zgodność
na \(A\) dają
\(S(v)=\sum_jc_jS(a_j)=\sum_jc_jT(a_j)=T(v)\).
Dotyczy to również \(V=\{0_V\}\)
i pustego zbioru generatorów.
\end{proof}

\begin{przyklad}[Dwa operatory na wielomianach]
\label{prz:lin3-przesuniecia-wielomianow}
Niech \(V=K[X]\) oraz \(B=(1,X,X^2,\ldots)\).
Wybór \(T(X^n)=X^{n+1}\) dla \(n\in\mathbb N_0\)
wyznacza dokładnie jedno odwzorowanie liniowe,
a z jego wzoru wynika \(T(p)=Xp\).
Z kolei wybór \(R(1)=0\) oraz
\(R(X^{n+1})=X^n\) dla \(n\in\mathbb N_0\)
wyznacza liniowe odwzorowanie
\[
 R(p)=\frac{p(X)-p(0)}{X}.
\]
Prawa strona oznacza wielomian, ponieważ
licznik ma zerowy wyraz stały.
Dla \(p=\sum_{j=0}^{d}a_jX^j\)
obliczamy \(R(p)=\sum_{j=1}^{d}a_jX^{j-1}\).
W szczególności \(R(T(p))=p\), ale
\(T(R(p))=p-p(0)\); oba złożenia
nie muszą być równe.
\end{przyklad}

\section{Jądro, obraz i warunki iniektywności}
\label{sec:lin3-jadro-obraz}

\begin{definicja}[Jądro i obraz]
\label{def:lin3-jadro-obraz}
Niech \(T\colon V\to W\) będzie odwzorowaniem
liniowym przestrzeni wektorowych nad \(K\).
Przez \emph{jądro} i \emph{obraz} rozumiemy zbiory
\[
 \ker T=\{v\in V:T(v)=0_W\},\qquad
 \operatorname{im}T=\{T(v):v\in V\}\subseteq W.
\]
Jądro jest zbiorem wektorów tracących
informację pod działaniem \(T\),
a obraz zbiorem wektorów osiągalnych w \(W\).
\end{definicja}

\begin{twierdzenie}[Jądro i obraz są podprzestrzeniami]
\label{thm:lin3-jadro-obraz-podprzestrzenie}
Niech \(T\colon V\to W\) będzie odwzorowaniem
liniowym przestrzeni wektorowych nad \(K\).
Wtedy \(\ker T\) jest podprzestrzenią \(V\),
a \(\operatorname{im}T\) jest podprzestrzenią \(W\).
\end{twierdzenie}
\begin{proof}
Wektor \(0_V\) należy do \(\ker T\),
ponieważ \(T(0_V)=0_W\).
Dla \(v,u\in\ker T\) oraz \(a,b\in K\)
otrzymujemy \(T(av+bu)=aT(v)+bT(u)=0_W\).
Zatem \(av+bu\in\ker T\), co na mocy
kryterium podprzestrzeni pokazuje pierwszą tezę.

Obraz zawiera \(0_W=T(0_V)\).
Jeśli \(y_1=T(v_1)\), \(y_2=T(v_2)\)
należą do obrazu, to
\(ay_1+by_2=T(av_1+bv_2)\)
również do niego należy.
Obraz jest więc podprzestrzenią \(W\).
\end{proof}

\begin{stwierdzenie}[Obrazy generatorów i przeciwobrazy podprzestrzeni]
\label{prop:lin3-generatory-przeciwobrazy}
Niech \(V,W\) będą przestrzeniami wektorowymi
nad \(K\), \(T\colon V\to W\) odwzorowaniem liniowym,
a \(A\subseteq V\) dowolnym zbiorem.
Wtedy
\[
 T(\operatorname{span}(A))
      =\operatorname{span}(T(A)).
\]
Jeżeli \(Z\subseteq W\) jest podprzestrzenią,
to \(T^{-1}(Z)=\{v\in V:T(v)\in Z\}\)
jest podprzestrzenią \(V\).
\end{stwierdzenie}
\begin{proof}
Dla \(v=\sum_{j=1}^m a_ja_j\), \(a_j\in A\),
liniowość daje \(T(v)=\sum_{j=1}^m a_jT(a_j)\),
a więc zawieranie z lewej do prawej.
Każdy wektor prawej strony ma postać
\(\sum_{j=1}^m a_jT(a_j)\), która jest obrazem
wektora \(\sum_{j=1}^m a_ja_j\in\operatorname{span}(A)\).
Stąd także zawieranie odwrotne.

Ponieważ \(T(0_V)=0_W\in Z\),
zbiór \(T^{-1}(Z)\) nie jest pusty.
Dla \(v,u\in T^{-1}(Z)\) i \(a,b\in K\)
mamy \(T(av+bu)=aT(v)+bT(u)\in Z\),
bo \(Z\) jest podprzestrzenią.
Spełnione jest więc kryterium podprzestrzeni.
\end{proof}

\begin{twierdzenie}[Jądro a iniektywność, obraz a surjektywność]
\label{thm:lin3-iniekcja-surjekcja}
Niech \(T\colon V\to W\) będzie odwzorowaniem
liniowym przestrzeni wektorowych nad \(K\).
Wtedy:
\begin{enumerate}[label=(\roman*)]
\item \(T\) jest iniektywne wtedy i tylko wtedy,
gdy \(\ker T=\{0_V\}\);
\item \(T\) jest surjektywne wtedy i tylko wtedy,
gdy \(\operatorname{im}T=W\);
\item \(T\) jest bijektywne wtedy i tylko wtedy,
gdy \(\ker T=\{0_V\}\) oraz \(\operatorname{im}T=W\).
\end{enumerate}
\end{twierdzenie}
\begin{proof}
Jeśli \(T\) jest iniektywne i \(v\in\ker T\),
to \(T(v)=0_W=T(0_V)\), a więc \(v=0_V\).
Odwrotnie, zakładając \(\ker T=\{0_V\}\),
z równości \(T(v)=T(u)\) dostajemy
\(T(v-u)=0_W\), czyli \(v-u=0_V\),
a zatem \(v=u\). Druga równoważność
wynika bezpośrednio z definicji surjektywności
i obrazu. Trzecia jest połączeniem pierwszych dwóch.
\end{proof}

\begin{przyklad}[Jądra i obrazy na wielomianach]
\label{prz:lin3-jadro-przesuniecia}
Dla operatorów \(T,R\colon K[X]\to K[X]\)
z Przykładu~\ref{prz:lin3-przesuniecia-wielomianow}
mamy
\[
 \ker T=\{0\},\quad
 \operatorname{im}T=XK[X],\qquad
 \ker R=K,\quad
 \operatorname{im}R=K[X].
\]
Pierwsze dwie równości wynikają z tego,
że iloczyn \(Xp\) jest zerowy tylko dla \(p=0\),
a wielomiany postaci \(Xp\) są dokładnie
wielomianami o zerowym wyrazie stałym.
W przypadku \(R\), równość \(R(p)=0\)
jest równoważna \(p=p(0)\), czyli stałości \(p\).
Dla dowolnego \(q\in K[X]\) mamy
\(R(Xq)=q\), więc \(R\) jest surjektywne.
Wynik pozostaje prawdziwy w każdej
charakterystyce ciała \(K\).
\end{przyklad}

\section{Wymiar jądra i obrazu}
\label{sec:lin3-twierdzenie-rzad}

\begin{definicja}[Rząd i wymiar jądra]
\label{def:lin3-rzad-nullity}
Niech \(T\colon V\to W\) będzie liniowe.
Przyjmując aksjomat wyboru dla dowolnych
przestrzeni wektorowych, definiujemy
\[
 \operatorname{rk}T:=\dim_K(\operatorname{im}T),
 \qquad
 \operatorname{null}T:=\dim_K(\ker T).
\]
Jeżeli \(V\) jest skończenie wymiarowa,
obie wielkości są liczbami naturalnymi
i można je określić bez aksjomatu wyboru.
\end{definicja}

\begin{twierdzenie}[Twierdzenie o wymiarze jądra i obrazu]
\label{thm:lin3-rzad-nullity}
Niech \(V,W\) będą przestrzeniami wektorowymi
nad \(K\), a \(T\colon V\to W\) odwzorowaniem liniowym.
Przy założeniu aksjomatu wyboru zachodzi
równość liczb kardynalnych
\[
 \dim_K V=\dim_K(\ker T)+\dim_K(\operatorname{im}T).
\]
Jeżeli \(V\) jest skończenie wymiarowa,
to \(\ker T\) i \(\operatorname{im}T\)
są skończenie wymiarowe, a ta równość
zachodzi również \emph{bez aksjomatu wyboru}.
W szczególności w tym przypadku
\(\operatorname{rk}T=\dim_KV-\operatorname{null}T\).
\end{twierdzenie}
\begin{proof}
Najpierw przyjmijmy aksjomat wyboru.
Wybierzmy bazę \(C\) jądra \(\ker T\)
i rozszerzmy ją do bazy \(B=C\sqcup D\)
całej przestrzeni \(V\).
Pokażemy, że rodzina \(T(D)=(T(d))_{d\in D}\)
jest bazą \(\operatorname{im}T\).

Dla \(y\in\operatorname{im}T\) istnieje \(v\in V\)
takie, że \(y=T(v)\). Rozwijając \(v\) w bazie \(B\),
otrzymujemy skończoną sumę
\[
 v=\sum_{c\in F}a_cc+\sum_{d\in G}b_dd,
 \qquad F\subseteq C,\ G\subseteq D.
\]
Ponieważ \(T(c)=0_W\) dla \(c\in C\),
mamy \(y=\sum_{d\in G}b_dT(d)\).
Rodzina \(T(D)\) rozpina więc obraz.

Aby sprawdzić warunek niezależności liniowej,
niech \(G\subseteq D\) będzie skończony i
\(\sum_{d\in G}b_dT(d)=0_W\).
Wektor \(u:=\sum_{d\in G}b_dd\) należy
wówczas do \(\ker T=\operatorname{span}(C)\).
Możemy zatem zapisać \(u=\sum_{c\in F}a_cc\)
dla pewnego skończonego \(F\subseteq C\).
Odejmując obie kombinacje otrzymujemy
\[
 \sum_{d\in G}b_dd-\sum_{c\in F}a_cc=0_V.
\]
Ponieważ \(B=C\sqcup D\) jest bazą,
wszystkie współczynniki \(b_d\) i \(a_c\)
muszą być zerowe. Dowodzi to niezależności \(T(D)\),
a zarazem zapewnia, że różne indeksy \(d\)
odpowiadają różnym wektorom \(T(d)\).

Mamy więc \(\dim_K(\operatorname{im}T)=|D|\),
\(\dim_K(\ker T)=|C|\) oraz
\(\dim_KV=|C|+|D|\).
Otrzymujemy wzór, bez odejmowania
nieskończonych liczb kardynalnych.

Jeśli \(V\) ma skończoną bazę, to
\(\ker T\) ma skończoną bazę \(C\)
na mocy wyników skończenie wymiarowych,
którą rozszerzamy skończoną procedurą
do bazy \(C\sqcup D\) przestrzeni \(V\).
Cały argument powyżej od tej chwili
stosuje się dosłownie, ale liczności
\(|C|\) i \(|D|\) są liczbami naturalnymi.
Nie użyto w nim aksjomatu wyboru.
\end{proof}

\begin{przyklad}[Wymiar jądra i obrazu konkretnego odwzorowania]
\label{prz:lin3-rzad-r3}
Niech \(T\colon\mathbb R^3\to\mathbb R^2\)
będzie dane wzorem
\[
 T(x,y,z)=(x+y,y+z).
\]
Jest liniowe, bo obie współrzędne są liniowymi
kombinacjami współrzędnych wejściowych.
Warunek \(T(x,y,z)=(0,0)\) daje układ
\(x=-y\), \(z=-y\). Zatem
\(\ker T=\operatorname{span}\{(-1,1,-1)\}\),
więc \(\operatorname{null}T=1\).
Z drugiej strony \(T(1,0,0)=(1,0)\)
i \(T(0,0,1)=(0,1)\), zatem
\(\operatorname{im}T=\mathbb R^2\),
czyli \(\operatorname{rk}T=2\).
Sprawdzamy wzór \(3=1+2\).
W szczególności \(T\) jest surjektywne,
lecz nie jest iniektywne.
\end{przyklad}

\begin{uwaga}[Wymiary nieskończone]
\label{rem:lin3-rzad-nieskonczony}
W przypadku nieskończenie wymiarowym
z równości
\(\dim_KV=\dim_K(\ker T)+\operatorname{rk}T\)
nie wolno odzyskiwać rzędu przez zwykłe
odejmowanie liczb kardynalnych.
Na przykład dla operatora \(T(p)=Xp\)
na \(K[X]\) mamy \(\ker T=\{0\}\)
i \(\operatorname{rk}T=\aleph_0\).
Dla operatora \(R\) z
Przykładu~\ref{prz:lin3-jadro-przesuniecia}
mamy \(\dim_K\ker R=1\)
i \(\operatorname{rk}R=\aleph_0\),
a mimo to oba wzory dają
\(\dim_KK[X]=\aleph_0\).
\end{uwaga}

\section{Pierwsze twierdzenie o izomorfizmie}
\label{sec:lin3-izomorfizm}

\begin{twierdzenie}[Pierwsze twierdzenie o izomorfizmie]
\label{thm:lin3-pierwsze-izo}
Niech \(V,W\) będą przestrzeniami wektorowymi
nad \(K\), a \(T\colon V\to W\)
odwzorowaniem liniowym. Wtedy
\[
 \widetilde T\colon V/\ker T
 \longrightarrow\operatorname{im}T,\qquad
 v+\ker T\longmapsto T(v)
\]
jest dobrze określonym izomorfizmem
przestrzeni wektorowych. Nie jest potrzebny
aksjomat wyboru.
\end{twierdzenie}
\begin{proof}
Najpierw sprawdzamy dobrą określoność.
Jeśli \(v+\ker T=u+\ker T\), to
z Lematu~\ref{lem:lin2-warstwy}
wynika \(v-u\in\ker T\).
Stąd \(T(v)-T(u)=T(v-u)=0_W\),
czyli \(T(v)=T(u)\).
Zatem wzór nie zależy od przedstawiciela warstwy.

Dla \(a,b\in K\) i \(v,u\in V\) otrzymujemy
\[
 \begin{aligned}
 \widetilde T\bigl(a(v+\ker T)+b(u+\ker T)\bigr)
  &=\widetilde T\bigl((av+bu)+\ker T\bigr)\\
  &=T(av+bu)=aT(v)+bT(u).
 \end{aligned}
\]
Jest to liniowość \(\widetilde T\).
Jeśli \(\widetilde T(v+\ker T)=0_W\),
to \(T(v)=0_W\), a więc \(v\in\ker T\).
Wówczas \(v+\ker T=\ker T=0_{V/\ker T}\).
Otrzymujemy iniektywność.
Wreszcie dla każdego \(y\in\operatorname{im}T\)
istnieje \(v\in V\) z \(y=T(v)\);
wtedy \(y=\widetilde T(v+\ker T)\).
Dowodzi to surjektywności.

Izomorfizm jest zatem określony wyłącznie
przez \(T\), bez wybierania dopełnienia
algebraicznego jądra ani bazy.
\end{proof}

\begin{figure}[!htbp]
\centering
\begin{tikzpicture}[>=Latex,every node/.style={font=\small}]
 \node (V) at (0,1.9) {\(V\)};
 \node (W) at (6,1.9) {\(\operatorname{im}T\)};
 \node (Q) at (3,0) {\(V/\ker T\)};
 \draw[->,thick] (V)--node[above]{\(T\)} (W);
 \draw[->,thick] (V)--node[below left]{\(\pi_{\ker T}\)} (Q);
 \draw[->,thick] (Q)--node[below right]{\(\widetilde T\)} (W);
\end{tikzpicture}
\caption{Rozkład odwzorowania \(T\) na
odwzorowanie ilorazowe i izomorfizm na obraz:
\(T=\widetilde T\circ\pi_{\ker T}\).}
\label{fig:lin3-izo}
\end{figure}

\begin{stwierdzenie}[Kiedy odwzorowanie przechodzi do ilorazu?]
\label{prop:lin3-faktoryzacja-iloraz}
Niech \(V,W\) będą przestrzeniami wektorowymi
nad \(K\), niech \(U\subseteq V\) będzie
podprzestrzenią, a \(T\colon V\to W\) odwzorowaniem
liniowym. Istnieje liniowe odwzorowanie
\(S\colon V/U\to W\) z \(T=S\circ\pi_U\)
wtedy i tylko wtedy, gdy \(U\subseteq\ker T\).
Jeśli istnieje, jest jednoznaczne i spełnia
\(S(v+U)=T(v)\).
\end{stwierdzenie}
\begin{proof}
Jeśli \(T=S\circ\pi_U\), to dla \(u\in U\)
mamy \(\pi_U(u)=0_{V/U}\), stąd
\(T(u)=S(0_{V/U})=0_W\). Zatem
\(U\subseteq\ker T\).
Odwrotnie, jeśli \(U\subseteq\ker T\), to
\(T\) zeruje się na \(U\);
Twierdzenie~\ref{thm:lin2-wlasnosc-ilorazu}
daje istnienie i jednoznaczność \(S\)
oraz wskazany wzór.
\end{proof}

\begin{przyklad}[Wartość wielomianu w zerze]
\label{prz:lin3-ewaluacja}
Niech \(T\colon K[X]\to K\) będzie
dane przez \(T(p)=p(0)\).
Jest liniowe, bo
\((ap+bq)(0)=ap(0)+bq(0)\).
Wszystkie wielomiany o zerowym wyrazie
stałym są podzielne przez \(X\), zatem
\(\ker T=XK[X]\).
Odwzorowanie jest surjektywne:
dla \(c\in K\) wystarczy wziąć wielomian stały \(c\).
Pierwsze twierdzenie o izomorfizmie daje
\[
 K[X]/XK[X]\cong K,\qquad
 p+XK[X]\longmapsto p(0).
\]
Otrzymujemy jednoznaczne utożsamienie warstw
z wyrazami stałymi, bez wybierania
reprezentantów innych niż te użyte we wzorze.
\end{przyklad}

\section{Składanie i odwracalność odwzorowań liniowych}
\label{sec:lin3-skladanie}

\begin{stwierdzenie}[Składanie odwzorowań liniowych]
\label{prop:lin3-skladanie}
Niech \(U,V,W\) będą przestrzeniami
wektorowymi nad \(K\), a \(T\colon U\to V\),
\(S\colon V\to W\) odwzorowaniami liniowymi.
Wtedy złożenie \(S\circ T\colon U\to W\) jest liniowe,
\[
 \ker T\subseteq\ker(S\circ T),\qquad
 \operatorname{im}(S\circ T)\subseteq\operatorname{im}S.
\]
Ponadto odwzorowanie identycznościowe
\(\operatorname{id}_V\colon V\to V\)
jest liniowe i jest elementem neutralnym
składania tam, gdzie to składanie jest określone.
\end{stwierdzenie}
\begin{proof}
Dla \(u_1,u_2\in U\) i \(a,b\in K\)
mamy
\[
 (S\circ T)(au_1+bu_2)
 =S(aT(u_1)+bT(u_2))
 =aS(T(u_1))+bS(T(u_2)).
\]
Dowodzi to liniowości. Jeśli \(u\in\ker T\),
to \((S\circ T)(u)=S(0_V)=0_W\),
czyli \(u\in\ker(S\circ T)\).
Każdy wektor obrazu złożenia ma postać
\(S(T(u))\), więc należy do \(\operatorname{im}S\).
Dla odwzorowania identycznościowego
wszystkie warunki wynikają z
\(\operatorname{id}_V(av+bw)=av+bw
 =a\operatorname{id}_V(v)+b\operatorname{id}_V(w)\).
Równości \(T\circ\operatorname{id}_U=T\)
oraz \(\operatorname{id}_V\circ T=T\)
zachodzą punktowo.
\end{proof}

\begin{twierdzenie}[Izomorfizmy i odwzorowania odwrotne]
\label{thm:lin3-odwrotnosc}
Niech \(T\colon V\to W\) będzie odwzorowaniem
liniowym przestrzeni wektorowych nad \(K\).
Wówczas \(T\) jest bijektywne wtedy i tylko wtedy,
gdy istnieje liniowe \(S\colon W\to V\)
takie, że \(S\circ T=\operatorname{id}_V\)
i \(T\circ S=\operatorname{id}_W\).
W takim przypadku \(S=T^{-1}\) i \(T\)
nazywamy \emph{izomorfizmem}.
\end{twierdzenie}
\begin{proof}
Jeżeli \(S\) spełnia obie równości,
to z \(T(v)=T(u)\) wynika
\(v=S(T(v))=S(T(u))=u\), czyli \(T\)
jest iniektywne. Każdy \(w\in W\)
ma postać \(w=T(S(w))\), więc \(T\) jest surjektywne.

Odwrotnie, dla bijektywnego \(T\)
istnieje odwrotne odwzorowanie zbiorów
\(T^{-1}\colon W\to V\).
Dla \(w_1,w_2\in W\) i \(a,b\in K\)
niech \(v_i=T^{-1}(w_i)\).
Wówczas
\(T(av_1+bv_2)=aw_1+bw_2\).
Z iniektywności \(T\) otrzymujemy
\[
 T^{-1}(aw_1+bw_2)
 =aT^{-1}(w_1)+bT^{-1}(w_2).
\]
A więc odwrotność jest liniowa
i spełnia obie wymagane równości.
\end{proof}

\begin{stwierdzenie}[Izomorfizm zachowuje bazę i wymiar]
\label{prop:lin3-izomorfizm-baza}
Niech \(T\colon V\to W\) będzie izomorfizmem
przestrzeni wektorowych nad \(K\).
Jeżeli \(B\) jest bazą Hamela \(V\),
to \(T(B)\) jest bazą Hamela \(W\).
W szczególności przy aksjomacie wyboru
\(\dim_KV=\dim_KW\).
Jeżeli \(V,W\) są skończenie wymiarowe,
równość wymiarów nie wymaga aksjomatu wyboru.
\end{stwierdzenie}
\begin{proof}
Ponieważ \(T\) jest surjektywne,
\[
 W=T(V)=T(\operatorname{span}(B))
   =\operatorname{span}(T(B))
\]
na mocy Stwierdzenia~\ref{prop:lin3-generatory-przeciwobrazy}.
Jeżeli skończona kombinacja różnych elementów
\(T(b_1),\ldots,T(b_m)\), \(b_j\in B\),
jest zerowa, to
\(T(\sum_{j=1}^{m}a_jb_j)=0_W\).
Iniektywność \(T\) daje
\(\sum_{j=1}^{m}a_jb_j=0_V\),
a niezależność \(B\) wymusza \(a_j=0_K\)
dla wszystkich \(j\).
Odwzorowanie \(b\mapsto T(b)\) jest iniektywne
na \(B\), więc zachowuje też moc rodziny bazowej.
Przy skończonych bazach jest to równość liczb
naturalnych i nie wymaga wyboru.
\end{proof}

\begin{wniosek}[Kryterium bijektywności w równych wymiarach]
\label{cor:lin3-rownymiary}
Niech \(V,W\) będą skończenie wymiarowymi
przestrzeniami wektorowymi nad \(K\)
o tym samym wymiarze \(n\in\mathbb N_0\).
Wtedy dla każdego liniowego \(T\colon V\to W\)
następujące warunki są równoważne:
\(T\) jest iniektywne, \(T\) jest surjektywne,
\(T\) jest izomorfizmem.
\end{wniosek}
\begin{proof}
Z Twierdzenia~\ref{thm:lin3-rzad-nullity}
mamy \(n=\dim_K\ker T+\dim_K\operatorname{im}T\).
Jeżeli \(T\) jest iniektywne, to
\(\dim_K\ker T=0\), zatem obraz ma wymiar \(n\).
Jako podprzestrzeń \(W\) o tym samym wymiarze,
obraz musi być równy \(W\)
(Twierdzenie~\ref{thm:lin-wymiar-podprzestrzeni}).
Wtedy \(T\) jest surjektywne.
Jeżeli \(T\) jest surjektywne, to
\(\dim_K\operatorname{im}T=n\),
a więc wymiar jądra jest zerowy.
Przestrzeń o wymiarze zero jest zerowa,
więc \(T\) jest iniektywne.
Ostatnia równoważność wynika z definicji
izomorfizmu i Twierdzenia~\ref{thm:lin3-odwrotnosc}.
Przy \(n=0\) przestrzenie \(V,W\) są zerowe
i jedyne odwzorowanie liniowe jest izomorfizmem.
\end{proof}

\begin{przyklad}[Dlaczego wymiar skończony jest istotny?]
\label{prz:lin3-infty-nie-rownowazne}
Operatory \(T,R\colon K[X]\to K[X]\)
z Przykładu~\ref{prz:lin3-przesuniecia-wielomianow}
spełniają \(R\circ T=\operatorname{id}_{K[X]}\)
oraz \(T\circ R(p)=p-p(0)\).
Zatem \(T\) jest iniektywne, ale nie jest
surjektywne: żaden wielomian \(Xp\)
nie jest równy \(1\).
Natomiast \(R\) jest surjektywne,
ale nie jest iniektywne, bo \(R(1)=0\).
Mimo to dziedzina i przeciwdziedzina
obu odwzorowań mają ten sam nieskończony wymiar
\(\aleph_0\). Wniosek~\ref{cor:lin3-rownymiary}
nie przenosi się bez dodatkowych założeń
na przestrzenie nieskończenie wymiarowe.
\end{przyklad}

\section{Przestrzenie odwzorowań liniowych}
\label{sec:lin3-hom}

\begin{definicja}[Przestrzenie \(\operatorname{Hom}\) i \(\operatorname{End}\)]
\label{def:lin3-hom-end}
Dla przestrzeni wektorowych \(V,W\) nad \(K\)
oznaczamy przez
\[
 \operatorname{Hom}_K(V,W)
  :=\{T\colon V\to W:T\text{ jest liniowe}\}
\]
zbiór wszystkich odwzorowań liniowych.
Przez \(\operatorname{End}_K(V)=\operatorname{Hom}_K(V,V)\)
rozumiemy zbiór \emph{endomorfizmów} \(V\).
Na \(\operatorname{Hom}_K(V,W)\) określamy działania
\[
 (S+T)(v)=S(v)+T(v),\qquad
 (aT)(v)=aT(v).
\]
\end{definicja}

\begin{twierdzenie}[Przestrzeń liniowa odwzorowań]
\label{thm:lin3-hom-przestrzen}
Niech \(V,W\) będą przestrzeniami wektorowymi nad \(K\).
Zbiór \(\operatorname{Hom}_K(V,W)\)
z działaniami punktowymi jest przestrzenią wektorową nad \(K\).
Jeżeli \(\dim_KV=n\) i \(\dim_KW=m\)
są skończone, to
\[
 \dim_K\operatorname{Hom}_K(V,W)=mn.
\]
Druga część twierdzenia nie wymaga aksjomatu wyboru.
\end{twierdzenie}
\begin{proof}
Najpierw \(\operatorname{Hom}_K(V,W)\)
zawiera odwzorowanie zerowe.
Jeżeli \(S,T\) są liniowe i \(a,b\in K\),
to dla dowolnych \(v,u\in V\) i \(c,d\in K\)
\[
 \begin{aligned}
 (aS+bT)(cv+du)
  &=aS(cv+du)+bT(cv+du)\\
  &=c(aS+bT)(v)+d(aS+bT)(u).
 \end{aligned}
\]
Suma i iloczyn przez skalar pozostają więc liniowe.
Pozostałe aksjomaty przestrzeni wektorowej
wynikają z odpowiednich aksjomatów \(W\)
po podstawieniu dowolnego \(v\in V\).

Niech teraz \(B=(v_1,\ldots,v_n)\)
i \(C=(w_1,\ldots,w_m)\) będą bazami odpowiednio
\(V\) i \(W\).
Dla każdej pary \(1\leq i\leq n\),
\(1\leq j\leq m\) Twierdzenie~\ref{thm:lin3-przedluzenie-z-bazy}
daje jednoznaczne liniowe \(E_{ji}\colon V\to W\),
określone na bazie przez
\[
 E_{ji}(v_k)=
 \begin{cases}
   w_j,&k=i,\\
   0_W,&k\neq i.
 \end{cases}
\]
Pokażemy, że rodzina \((E_{ji})\) jest bazą
\(\operatorname{Hom}_K(V,W)\).

Dla dowolnego liniowego \(T\colon V\to W\)
rozwińmy \(T(v_i)=\sum_{j=1}^{m}a_{ji}w_j\).
Odwzorowania \(T\) oraz
\(\sum_{i=1}^{n}\sum_{j=1}^{m}a_{ji}E_{ji}\)
zgadzają się na każdym \(v_k\);
na mocy Wniosku~\ref{cor:lin3-zgodnosc-na-generatorach}
są więc równe. Dowodzi to rozpinania.

Jeżeli \(\sum_{i,j}c_{ji}E_{ji}=0\),
to po zastosowaniu obu stron do \(v_k\)
otrzymujemy \(\sum_{j=1}^{m}c_{jk}w_j=0_W\).
Niezależność bazy \(C\) daje
\(c_{jk}=0_K\) dla każdego \(j,k\).
Jest to warunek niezależności liniowej
rodziny \((E_{ji})\). Jej liczba elementów
wynosi \(mn\), co dowodzi wzoru.

Przy \(n=0\) istnieje tylko odwzorowanie zerowe
z \(V=\{0_V\}\) do \(W\); przy \(m=0\)
istnieje tylko odwzorowanie zerowe do \(W=\{0_W\}\).
W obu przypadkach \(\operatorname{Hom}_K(V,W)\)
jest przestrzenią zerową, a wzór daje \(mn=0\).
\end{proof}

\begin{stwierdzenie}[Algebra endomorfizmów]
\label{prop:lin3-end-algebra}
Niech \(V\) będzie przestrzenią wektorową nad \(K\).
Przestrzeń \(\operatorname{End}_K(V)\)
z dodawaniem i mnożeniem przez skalary
wykonywanymi punktowo oraz z mnożeniem
\(S\cdot T:=S\circ T\)
jest łączną algebrą nad \(K\) z jedynką
\(\operatorname{id}_V\).
W ogólności mnożenie nie jest przemienne.
\end{stwierdzenie}
\begin{proof}
Dzięki Stwierdzeniu~\ref{prop:lin3-skladanie}
złożenie endomorfizmów jest endomorfizmem.
Łączność wynika z łączności składania funkcji.
Dla \(S,T,R\in\operatorname{End}_K(V)\),
\(a\in K\) i \(v\in V\) mamy
\[
 \begin{aligned}
 ((S+T)\circ R)(v)&=S(R(v))+T(R(v)),\\
 (S\circ(T+R))(v)&=S(T(v)+R(v))
                  =S(T(v))+S(R(v)),\\
 ((aS)\circ T)(v)&=aS(T(v))
                  =(S\circ(aT))(v).
 \end{aligned}
\]
Są to rozdzielności i zgodność
mnożenia z mnożeniem przez skalary.
Jedynką jest odwzorowanie identycznościowe.

Aby zobaczyć, że nie musi zachodzić
przemienność, rozważmy na \(K^2\)
endomorfizmy \(P(x,y)=(x,0)\)
oraz \(Q(x,y)=(y,x)\).
Mamy \((P\circ Q)(1,0)=(0,0)\),
natomiast \((Q\circ P)(1,0)=(0,1)\).
Wobec \(1_K\neq0_K\) oba złożenia są różne.
\end{proof}

\section{Projekcje i sumy proste}
\label{sec:lin3-projekcje}

\begin{definicja}[Projekcja algebraiczna]
\label{def:lin3-projekcja}
Niech \(V\) będzie przestrzenią wektorową nad \(K\).
Endomorfizm \(P\colon V\to V\) nazywamy
\emph{projekcją algebraiczną}, jeśli
\(P^2=P\), gdzie \(P^2=P\circ P\).
\end{definicja}

\begin{twierdzenie}[Projekcje i rozkłady na sumę prostą]
\label{thm:lin3-projekcja-suma}
Niech \(V\) będzie przestrzenią wektorową nad \(K\)
i \(P\colon V\to V\) endomorfizmem.
Wtedy \(P^2=P\) wtedy i tylko wtedy, gdy
\[
 V=\operatorname{im}P\oplus\ker P
\]
oraz \(P\) jest identycznością na
\(\operatorname{im}P\).

Ponadto jeśli \(V=U\oplus W\)
dla podprzestrzeni \(U,W\subseteq V\),
to istnieje dokładnie jedna projekcja liniowa
\(P\colon V\to V\) spełniająca
\(P(u)=u\) dla \(u\in U\) i \(P(w)=0_V\)
dla \(w\in W\). Wówczas
\(\operatorname{im}P=U\) i \(\ker P=W\).
\end{twierdzenie}
\begin{proof}
Załóżmy \(P^2=P\). Każdy \(v\in V\)
ma rozkład
\[
 v=P(v)+(v-P(v)).
\]
Pierwszy składnik należy do obrazu,
a drugi do jądra, ponieważ
\(P(v-P(v))=P(v)-P^2(v)=0_V\).
Zatem \(V=\operatorname{im}P+\ker P\).
Jeżeli \(x\in\operatorname{im}P\cap\ker P\),
to \(x=P(y)\) dla pewnego \(y\in V\),
stąd \(P(x)=P^2(y)=P(y)=x\),
ale zarazem \(P(x)=0_V\).
Wobec tego \(x=0_V\), a suma jest prosta.
Rachunek \(P(P(y))=P(y)\) pokazuje również,
że \(P\) jest identycznością na obrazie.

W odwrotnym kierunku załóżmy,
że \(V=\operatorname{im}P\oplus\ker P\)
oraz \(P(x)=x\) dla \(x\in\operatorname{im}P\).
Dla każdego \(v\in V\) wektor \(P(v)\)
należy do obrazu, więc \(P^2(v)=P(v)\);
mamy zatem \(P^2=P\).

Niech teraz \(V=U\oplus W\).
Z Twierdzenia~\ref{thm:lin2-kryteria-sumy-prostej}
każdy \(v\in V\) ma dokładnie jeden rozkład
\(v=u+w\), \(u\in U,w\in W\).
Definiujemy \(P(v)=u\).
Dla \(v_i=u_i+w_i\) i \(a,b\in K\)
wektor \(av_1+bv_2\) ma rozkład
\((au_1+bu_2)+(aw_1+bw_2)\);
jednoznaczność daje
\(P(av_1+bv_2)=au_1+bu_2
 =aP(v_1)+bP(v_2)\).
Odwzorowanie jest więc liniowe.
Ponieważ \(P(u+w)=u\), mamy
\(P^2=P\), \(\operatorname{im}P=U\)
i \(\ker P=W\).
Każde inne odwzorowanie liniowe o wskazanych
wartościach na \(U\) i \(W\) musiałoby spełniać
\(Q(u+w)=Q(u)+Q(w)=u=P(u+w)\).
Dowodzi to jednoznaczności.
\end{proof}

\begin{przyklad}[Projekcja na oś]
\label{prz:lin3-projekcja-os}
W \(\mathbb R^2\) przyjmijmy
\(U=\mathbb R\times\{0\}\) oraz
\(W=\{0\}\times\mathbb R\).
Projekcja wzdłuż \(W\) na \(U\) jest dana wzorem
\[
 P(x,y)=(x,0).
\]
Wtedy \(P^2=P\), \(\operatorname{im}P=U\),
\(\ker P=W\) i \(\mathbb R^2=U\oplus W\).
Jeśli zamiast \(W\) wybierzemy
\(\operatorname{span}\{(1,1)\}\),
to otrzymamy inne dopełnienie \(U\)
i inną projekcję \(P'(x,y)=(x-y,0)\).
Potwierdza to niejednoznaczność dopełnień
z Przykładu~\ref{prz:lin2-dopelnienia}.
\end{przyklad}

\section{Podsumowanie}
\label{sec:lin3-podsumowanie}

Odwzorowanie liniowe jest wyznaczone jednoznacznie
przez wartości na ustalonej bazie Hamela.
Jego jądro i obraz są podprzestrzeniami,
a jądro charakteryzuje iniektywność.
W przypadku ogólnym przy aksjomacie wyboru
obowiązuje kardynalny wzór
\[
 \dim_KV=\dim_K\ker T+\dim_K\operatorname{im}T,
\]
który w skończonym wymiarze nie wymaga tego aksjomatu.
Zarazem zawsze, bez aksjomatu wyboru, mamy
kanoniczny izomorfizm
\(V/\ker T\cong\operatorname{im}T\).
Odwzorowania liniowe tworzą przestrzeń
\(\operatorname{Hom}_K(V,W)\), a endomorfizmy
można dodatkowo składać.
Projekcje algebraiczne opisują dokładnie
rozkłady przestrzeni na sumy proste
wzdłuż wskazanych dopełnień.
